\documentclass[lineno]{biometrika}

\usepackage{amsmath}
%\usepackage{graphics}

%% Please use the following statements for
%% managing the text and math fonts for your papers:
%\usepackage{times}
%\usepackage[cmbold]{mathtime}
%\usepackage{bm}

\usepackage{newtxtext}
\usepackage[subscriptcorrection]{newtxmath}

%\usepackage{natbib}

\graphicspath{{./art/}}

\usepackage[plain,noend]{algorithm2e}

\makeatletter
\renewcommand{\algocf@captiontext}[2]{#1\algocf@typo. \AlCapFnt{}#2} % text of caption
\renewcommand{\AlTitleFnt}[1]{#1\unskip}% default definition
\def\@algocf@capt@plain{top}
\renewcommand{\algocf@makecaption}[2]{%
  \addtolength{\hsize}{\algomargin}%
  \sbox\@tempboxa{\algocf@captiontext{#1}{#2}}%
  \ifdim\wd\@tempboxa >\hsize%     % if caption is longer than a line
    \hskip .5\algomargin%
    \parbox[t]{\hsize}{\algocf@captiontext{#1}{#2}}% then caption is not centered
  \else%
    \global\@minipagefalse%
    \hbox to\hsize{\box\@tempboxa}% else caption is centered
  \fi%
  \addtolength{\hsize}{-\algomargin}%
}
\makeatother

%%% User-defined macros should be placed here, but keep them to a minimum.
\def\Bka{{\it Biometrika}}

\def\AIC{\textsc{aic}}
\def\T{{ \mathrm{\scriptscriptstyle T} }}
\def\v{{\varepsilon}}

%\addtolength\topmargin{35pt}
\DeclareMathOperator{\Thetabb}{\mathcal{C}}

%\usepackage{hyperref}

\begin{document}

\jname{Biometrika}
%% The year, volume, and number are determined on publication
\jyear{2025}
\jvol{112}
\jnum{1}
\cyear{2025}
%% The \doi{...} and \accessdate commands are used by the production team
%\doi{10.1093/biomet/asm023}
\accessdate{Advance Access publication on 14 February 2025}

%% These dates are usually set by the production team
\received{2 January 2024}
\revised{3 February 2025}

%% The left and right page headers are defined here:
\markboth{P. Fearnhead et~al.}{Biometrika style}

%% Here are the title, author names and addresses
\title{Some notes on \textit{Biometrika} style}

\author{P. FEARNHEAD}
\affil{Department of Mathematics and Statistics, Lancaster University,\\ Bailrigg, Lancaster LA1 4YF, U.K.
\email{editor@biometrikatrust.org}}

\author{A. C. DAVISON}
\affil{Institute of Mathematics, Ecole Polytechnique F\'ed\'erale de Lausanne,\\ Station 8, 1015 Lausanne, Switzerland \email{anthony.davison@epfl.ch}}


\author{R. E. GESSNER}
\affil{Department of Statistical Science, University College London,\\ Gower Street, London WC1E 6BT, U.K.
\email{editorial@biometrikatrust.org}}

\author{\and D. M. TITTERINGTON}
%\affil{Department of Statistics, University of Glasgow, Glasgow G12 8QQ, U.K. \email{mike@stats.gla.ac.uk}}

\maketitle

\begin{abstract}
There should be a single paragraph summary which should not contain formulae, symbols or acronyms, followed by 3--8 key words in alphabetical order.   The summary contains bibliographic references only if they are essential.  It should indicate results rather than describe the contents of the paper: for example, `A simulation study is performed' should be replaced by a more informative phrase such as `In a simulation our estimator had smaller mean square error than its main competitors.'
\end{abstract}

\begin{keywords}
Address; Appendix; Figure; Length; Reference; Style; Summary; Table.
\end{keywords}

\section{Introduction}

These notes are intended both as a guide to journal style and as an example file or template for authors intending to prepare a paper using  the \Bka\ \LaTeX\ class, which is encouraged.  It is not necessary for papers submitted to the journal to follow these notes, and a decision on a paper will not be affected by the level that it conforms to the style requirements. However papers that are accepted may require revisions to adhere to the style of the journal.

Where the primary contribution of a paper is a new method, code implementing the method should be made available, either as part of the Supplementary Material or via a public repository, e.g., on GitHub. If simulation results are central to evaluation of such methods, then details of how to reproduce those results should also be available. 

\Bka\  papers do not contain a `contents of the paper' paragraph.

\section{Address}

For each author please give one postal address, including a department, postcode and country, and one e-mail address; these should be the best permanent addresses current at time of publication. Acknowledgements to other institutions should be put with other acknowledgements at the end of the paper. 

\section{Scope and length}

There are no restrictions on the area of statistics considered by Biometrika. The three
primary criteria by which all papers are judged are:
(i) importance of the work to the understanding, practice or potential practice of statistics;
(ii) the degree of conceptual novelty and insight;
(iii) the conciseness and clarity with which the ideas are conveyed.

The novelty may take many forms, including: the development of new methodology
accompanied by an insightful analysis; papers where the innovation is around computational
efficiency, where this is important for the application of the statistical method; a penetrating
exposition of anomalous or unforeseen behaviour of mainstream inferential tools; original
formulations uncovering foundational structure with potential relevance to statistical
practice. Papers concerned purely with sampling properties of existing procedures or minor
developments thereof are typically unsuitable unless such properties reveal considerable
structural understanding. We do not publish purely applied work.

Biometrika publishes three types of paper: regular papers, synthesis papers (both normally
fewer than 20 pages) and miscellanea articles (max. 8 pages). Miscellanea papers are not
expected to have the breadth of contribution of regular papers but need to meet the same
standards in terms of importance of the work, novelty of the insights, and general quality.
Synthesis papers should either open up an emerging area outside of statistics to a statistical
audience, assuming that area will be important to statistics, or be a synthesis of two or more
areas of statistics that brings fundamental new insight.

\section{Use of generative AI and AI-assisted technologies}
Authors are permitted to use generative AI and AI-assisted technologies in the writing process. This must be done with oversight by authors, who should review and, where necessary, edit the output. The authors must take full responsibility for all content of the paper. In particular, AI and AI-assisted technologies do not qualify as authors, and the Journal will screen for them in author lists. Please see the COPE position statement on Authorship and AI for more details.
 
Authors must disclose their use of AI and AI-assisted technologies, both within the cover letter to the editor and within the manuscript. This should detail what AI-technology is used and for what reason. Where AI has been used, please also add a statement within the manuscript. This should appear before the Acknowledgement section, and be titled ``Declaration of the use of generative AI and AI-assisted technologies''. It should read:
 
``During the preparation of this work the author(s) used [NAME TOOL/SERVICE] in order to [REASON]. After using this tool/service the author(s) reviewed and edited the content as necessary and take(s) full responsibility for the content of the publication.''
 
This declaration does not apply to the use of basic tools for checking grammar, spelling etc. If there is nothing to declare there is no need for a statement.
 
Papers will be judged solely on their content, and the use of generative AI within the writing of the paper will not prejudice the review of the paper.



\section{Style}

\subsection{Sections, subsections and paragraphs}

If subsections are used to divide a section, no text should appear before the first subsection; all text should appear within numbered subsections.  Subsubsections are not used.

\subsection{Some conventions}

English spelling is used, with Oxford ``-ize'' endings.

Verbal phrases inside dashes, or in  bold type, should not be used  and phrases inside brackets should be used sparingly. Parenthetical remarks, like this subordinate clause, are placed between two commas. Direct quotations should be attributed. Footnotes should be avoided except for tables. \Bka\  allows enumerated lists but not bullet point lists. 


Abbreviations should be used sparingly, and only if they improve readability. 
Do not create abbreviations to describe methods. Thus `our method is more efficient than Wellner and Zhang's method' should replace `our new method is better than method WZ'.


Symbols comprising several letters such as \AIC\ or \textsc{ar}$(p)$ may be used as mathematical objects if previously defined.   In such cases small capital letters, for example the  \TeX\ syntax \verb+\textsc{aic}+ for \AIC,\ are used; consistency is best assured by defining a macro at the start of the \texttt{.tex}  file.

\subsection{English}

English sentences containing mathematical expressions or displayed formulae should be punctuated in the usual way: in particular please check carefully that all displayed expressions are correctly punctuated.  

Words in common terms such as central limit theorem or Brownian motion are capitalized only if they are derived from proper names: thus bootstrap,  lasso and mean square error rather than Bootstrap, Lasso and  Mean Square Error.
 

Two bugbears: the phrase `note that' can almost always be deleted, and the phrase `is given by' should be cut to `is' in a sentence such as `The average is given by $\bar X=n^{-1}(X_1+\cdots+X_n)$'.

\subsection{Mathematics}

 Equation numbers should be included only when equations are referred to; the
numbers must be placed on the right. Long or important mathematical, not verbal,
expressions should be displayed, i.e., shown on a separate line. Short formulae should be
left in the text to save space where possible, but must not be more than one line high and
not contain reduced-size type. For example $\frac{{\rm d}y}{{\rm d}x}$ must not be left in the text, but should be
written ${{\rm d}y}/{{\rm d}x}$ or it should be displayed. Likewise write $n^{1/2}$ not $n^{\frac12}$.   Also $\displaystyle{a\choose{b}}$ and suchlike expressions
 must not be left in the text. Equations
involving lengthy expressions should, where possible, be avoided by introducing suitable
notation.

Symbols should not start sentences.  Vectors are assumed to be column vectors, unless
explicitly transposed. The use of an apostrophe to denote matrix or vector transposition should be avoided; it is preferable to write $A^\T$, $a^\T$.   Capital script letters may be used sparingly, typically to denote sets,
but care should be taken as some are hard to distinguish.

For equations with multiple nested brackets, please ensure readability by using different pairs, e.g., $[\{(\ldots)\}]$ as appropriate.   Follow the usual conventions for $e$, exp, use of solidus, square root signs and so forth
as in a recent issue. The sign $\sqrt{\hphantom{a}}$ is not used, and  the sign $\surd$ is used
only sparingly; powers of complicated quantities should be represented as $(mnpq)^{a}$.

Multiple overbars such as $\bar{\tilde{x}}$ must be avoided, as must $\widehat{ab}, (\widehat{a+b}), \widehat{\rm var}, \overline{ab},(\overline{a+b})$. Subscripts and superscripts, and second-order sub- and
superscripts, should be aligned horizontally. Avoid sub- and superscripts of third, and greater, order.

Please use: var($x$) not var $x$ or Var($x$); cov not Cov; pr for probability not Pr or P;
tr not trace; $E(X)$ for expectation not $EX$ or ${\cal E}(X)$; $\log x$ not $\log_e x$ or $\ln x$; $r$th not $r$-th or $r^{\rm th}$.
Please avoid: `$\cdot$' or `.' for product; $a/bc$, which should be written $a/(bc)$ or $a(bc)^{-1}$. Use the form $x_1,\ldots , x_n$ not $x_1, x_2,\ldots x_n$ and $\sum^{n}_{i=1}$ not $\sum^n_1$. Zeros precede decimal points: 0.2 not .2.

The use of `$\cdots$' and `$\ldots$' is  $\ldots$ in lists, such as $y_1,\ldots,y_n$, and $\cdots$ between binary operators,  giving $y_1+\cdots+y_n$.
Ranges of integers are denoted $i=1,\ldots, n$, whereas $0\leq x\leq 1$ is used for ranges of real numbers.  %The support of density and other functions should always be given.

\Bka\ deprecates the appearance of words in displayed equations, which should be formatted as
\begin{equation}
\label{one}
\bar Y = n^{-1} \sum_{j=1}^n Y_j,\quad S^2 = \sum_{j=1}^n (Y_j-\bar Y)^2;
\end{equation}
note the punctuation and space between the expressions. Displays such as \eqref{one} should take no more space than necessary, being placed on a single line where possible. Indexed equations and similar quantities in text are formatted as
$y_j = x_j^\T\beta + \v_j\ (j=1,\ldots, n)$, and are displayed as
\[
y_{ij} = x_j^\T\beta_i + \v_{ij}\quad (i=1,\ldots, m;\ j=1,\ldots, n).
\]


\subsection{Figures}


\begin{figure}
% The arguments in the next line are {height}{optional width}{used only by OUP for typesetting} for figure empty box eg 
%\figurebox{20pc}{25pc}{}
%if actual size of graphics need plese see below command
%\figurebox{}{}{}[fig1]
%need to reducing the figure size use below command
%\figuresize{.8}%
\figuresize{.6}
\figurebox{20pc}{25pc}{}[fig1]
% note that files may not be rotated
\caption{A graph showing the truth (dot-dash), an estimate (dashes), another estimate (solid), and 95\% pointwise confidence limits (small dashes).}
\label{fig1}
\end{figure}

%Figures are a common source of delay during production, usually because elementary guidelines have not been respected.  General comments may be found in the document `RSSGraphs.pdf' enclosed with the \Bka\ formatting files, and more detail is given in standard references such as \nocite{Cleveland:1993,Cleveland:1994} or \nocite{Tufte:1983}.

All the elements of a graph, including axis labels, should be large enough to be read easily, so the graph should be given a shape that will use the page space well. The use of large symbols, such as $\times$, for points should be avoided. If both axes of a panel show the same quantities, the panel should usually be square.  

Check that all the axes are labelled correctly and include units of measurement.  Axis labels should have the format `Difference of loglikelihoods': only the initial letter of the first word is upper-case.



%Please submit figures in greyscale wherever possible.  \Bka\ publishes in colour where this is essential, but care should be taken to ensure that any colours chosen will be distinguishable on the cream paper used for the journal, i.e., yellow should be avoided.

Figures should be referred to consecutively by number. Use of the \LaTeX\ \verb+\label+ and \verb+\ref+ commands to refer to figures and tables helps to reduce errors and so is preferred.  Figure~\ref{fig1}  is a reference to a figure at the start of a sentence, whereas subsequent references are abbreviated, for example to Fig.~\ref{fig1}.


\begin{table}
	\tbl{Comparison of the four methods based on different\break sample sizes}
{	\begin{tabular}{crrrcrrr}
	%	\hline
		& \multicolumn{3}{c}{$n=500$} & &\multicolumn{3}{c}{$n=1000$} \\
		%\cmidrule{2-4} \cmidrule{6-8} 
		Method & Bias & SD & RMSE & &Bias & SD & RMSE \\ 
	%	\hline
		Smith et al. (2000) & $4$ & $10$ & $11$& & $-1$ & $4$ & $4.1$\\ 
		Zhou et al. (2005) & $-51$ & $10$ & $12$ &&  $-2$ & $5$ & $5.2$\\ 
		Small \& Chen (2010) & $-65$ & $10$ & $13$ && $-3$ & $6$ & $6.3$\\ 
		Our proposed method & $-76$ & $15$ & $14$ && $4$ & $7$ & $7.4$\\ 
	%	\hline
	\end{tabular}}
	\label{tablelabel}
	\begin{tabnote}
	SD, standard deviation; RMSE, root mean squared error. All values have been multiplied by $10^2$.
	\end{tabnote}
	\label{tab2}
\end{table}


\subsection{Tables}

Tables should be referred to consecutively by number. Table is not abbreviated to Tab.

Check that the arrangement makes effective use
of the \Bka\ page. Layouts that have to be turned sideways should be avoided
if possible. 
Often tables can be improved by multiplying all the entries by a power of ten, so that 0.002 and
0.02 become 2 and 20 respectively, for example; this will often both save space and convey the message of the table more effectively. 

Very often tables containing results of Monte Carlo simulations use more digits than can be justified by the size of the simulation, and space can be gained and clarity improved by appropriate rounding.
Standard errors or some other measure of precision should be given for Monte Carlo results.
Often it suffices to give a phrase such as `The largest standard error for the results in column 2 is 0.01.' in the caption to the table.



%\begin{table}
%\def~{\hphantom{0}}
%\tbl{Perceptions about racial groups in the U.S. population}{%
%\begin{tabular}{lcccc}
%%\\
% \\
%& 2000 Census& \multicolumn{3}{c}{Mean percent estimated for U.S.} \\
%& percent of  & \multicolumn{3}{c}{population } \\
%& U.S. population & White Rs & Black Rs & Hispanic Rs \\[5pt]
%White & 75 & 59 & 56 & 60 \\
%Black & 12 & 30 & 38 & 40 \\
%Asian & ~4 & 16 & 21 & 30 \\
%American Indian & ~1 & 13 & 17 & 23 \\
%More than two races & ~2 & 41 & 48 & 50 \\
%Hispanic & 13 & 23 & 27 & 42
%\end{tabular}}
%\label{tablelabel}
%\begin{tabnote}
%U.S., United States of America; R, respondent.
%\end{tabnote}
%\end{table}

\subsection{Captions to figures and tables}

The caption to a figure should contain descriptions of lines and symbols used, but these should not be duplicated in the text, which should give the interpretation of the figure.  Verify that the caption and the graph agree, that every line and symbol is described correctly and that all lines or symbols in the graph are described. Figure captions always end with a full stop. 

The caption to a table should be concise; interpretation of the table should be given in the text. Any abbreviations used in the body of a table should be explained in a footnote to it.



\subsection{References}

References in the text should follow the current style used in \Bka.\ It is preferable to use BibTeX
if possible, as in this guide.  In citing
references use `First author et al.' if there are three or more authors. The list of references at
the end should correspond to those in the text, and be in exactly current \Bka\
form.

References to books should be to the latest edition; a page, section or chapter number
is nearly always necessary. References to books of papers should include title of book,
editor(s), first and final page numbers of paper, where published and publisher.

Complete lists of authors and editors should be given; in exceptional cases they may be abbreviated at the discretion of the editor.

PhD theses, unpublished reports and articles can be referred to in the text, using a phrase such as `as shown in a 2009 Euphoric State University Department of Statistics PhD thesis by M.~Zapp' or `the proofs may be found in an unpublished 2003 technical report available from the first author', but should not be included in the References except where they have been accepted for publication, and unless they appear in a permanent repository such as arXiv; in this case the most recent version of the work is cited like a paper, e.g., \citet{Berrendero.etal:2015}.

 URLs should be presented using the \verb;\texttt{}; font in LaTeX. Avoid giving URLs of pages that are likely
to become obsolete quickly.  Technical details for published papers should be prepared as Supplementary Material, so that they remain permanently available. Likewise software should be submitted as Supplementary Material; it should be adequately documented, e.g., by including a README file to accompany code.

\citet{Cox:1972} is an example of an active citation, and an example of a passive citation is  \citep{Hear:Holm:Step:quan:2006}.  The abbreviations for their journals should be noted.

\subsection{Theorem-like environments}

\Bka\  allows enumerated lists but not bullet point lists.
In this subsection we illustrate the use of theorem-like environments.

\begin{definition}[Optional argument]
This is a definition.
\end{definition}

\begin{assumption}[Another optional argument]
\label{assumptionA}
This is an assumption.
\end{assumption}

\begin{proposition}
This is a proposition.
\end{proposition}

\begin{lemma}
\label{lemma1}
This lemma precedes a theorem.
\end{lemma}

\begin{proof}
This is a proof of Lemma~\ref{lemma1}.  Perhaps it should be placed in the Appendix.
\end{proof}

\begin{theorem}
\label{theorem1}
This is a theorem.
\end{theorem}

Some text before we give the proof.

\begin{proof}[of Theorem~\ref{theorem1}]
The proof should be here.
\end{proof}

\begin{example}
This is an example.
\end{example}

Some text before the next theorem.

\begin{theorem}[Optional argument]
Another important result.
\end{theorem}

\begin{corollary}
This is a corollary.
\end{corollary}

\begin{remark}
This is a remark.
\end{remark}

\begin{step}
This is a step.
\end{step}

\begin{condition}
This is a condition.
\end{condition}


\begin{property}
This is a property.
\end{property}

\begin{restrictions}
This is a restriction.
\end{restrictions}

\begin{algo}
A simple algorithm.
%\vspace*{-12pt}
\begin{tabbing}
   \qquad \enspace Set $s=0$\\
   \qquad \enspace For $i=1$ to $i=n$ \\
   \qquad \qquad Set $t=0$\\
   \qquad \qquad For $j=1$ to $j=i$ \\\
  \qquad \qquad\qquad  $t \leftarrow t + x_{ij}$ \\
\qquad \qquad $s \leftarrow s + t$ \\
\qquad \enspace Output $s$
\end{tabbing}
\end{algo}


%\begin{algorithm}[!h]
%\caption{A simple algorithm.} \label{al1}
%\end{algorithm}


\section{Discussion}

This is the concluding part of the paper.  It is only needed if it contains new material such as open questions and future research areas.
It  should not repeat the summary or reiterate the contents of the paper.


\section*{Declaration of the use of generative AI and AI-assisted technologies}

[ONLY INCLUDE IF GENERATIVE AI OR AI-ASSISTED TECHNOLOGIES HAVE BEEN USED IN PRODUCTION OF THE PAPER; SEE \S 4]. During the preparation of this work the author(s) used [NAME TOOL/SERVICE] in order to [REASON]. After using this tool/service the author(s) reviewed and edited the content as necessary and take(s) full responsibility for the content of the publication.


\section*{Acknowledgement}
Acknowledgements should appear after the body of the paper but before any appendices and be as brief as possible
subject to politeness. 
%Information, such as contract numbers, of no interest to readers, should be excluded.

\section*{Supplementary material}
\label{SM}
Further material such as technical details, extended proofs, code, for example in \texttt{R} \citep{R:2025}, or additional  simulations, figures and examples may appear online, and should be briefly mentioned as Supplementary Material where appropriate.  Please submit any such content as a PDF file along with your paper, entitled `Supplementary material for Title-of-paper'.
Include in the paper a section `Supplementary material',  with the sentence `The Supplementary Material includes $\ldots$', giving a brief indication of what is available. It should be possible to read and understand the paper without reading the Supplementary Material.



%Further instructions will be given when a paper is accepted.
%\vspace*{-10pt}

\appendix

\appendixone
\section*{Appendix 1}
\subsection*{General}

Any appendices appear after the acknowledgement but before the references, and have titles. If there is more than one appendix, then they are numbered,  as here Theorem~\ref{appthm1}.

\begin{theorem}\label{appthm1}
This is a rather dull theorem:
\begin{equation}
\label{A1}
a + b = b + a;
\end{equation}
a little equation like this should only be displayed and labelled if it is referred to  elsewhere. 
\end{theorem}

\begin{lemma}\label{lem:gdp-1}
If $\alpha_j > 2$, $\eta_j/\alpha_j = O(j^{-m})$ ($j=1, \ldots,
\infty$) and $m > 1/2$, then $P_{l} (\Thetabb)= 1$.
\end{lemma}



\vspace*{-10pt}
\appendixtwo
\section*{Appendix 2}
\subsection*{Technical details}

Often the appendices contain technical details of the main results.

\begin{theorem}
This is another theorem full of gory details.
\end{theorem}

\begin{lemma}\label{lem:gdp-2}
If $\delta > 2$, $\rho > 0$, $\alpha_j(\delta) = \delta^j$ and
$\eta_j(\rho) = \rho$ for $j = 1, \ldots, \infty$, then $P_{l}
(\Thetabb) = 1$, where $P_{l}$ has density $p_{\mathrm{mgdP}}$ in
(4) with hyperparameters $\alpha_j(\delta)$ and
$\eta_j(\rho)$ ($j = 1, \ldots, \infty$). Furthermore, given
$\epsilon > 0$, there exists a positive integer $k(p, \delta,
\epsilon)= O\{\log^{-1} \delta \log ({p}/{\epsilon^2})\}$ for every
$\Omega$ such that for all $r \geq k$, $\alpha_j(\delta)=\delta^j$,
$\eta_j(\rho)=\rho$ ($j=1, \ldots, r$) and $\Omega^{r} = \Lambda^{r}
{\Lambda^r}^{\T} + \Sigma$, we have that ${\rm pr}\{\Omega^{r} \mid
d_{\infty}(\Omega, \Omega^{r}) < \epsilon\} > 1 - \epsilon$ where
$d_{\infty}(A, B) = {\max}_{1 \leq i,j \leq p}|a_{ij} - b_{ij}|$.
\end{lemma}

\vspace*{-10pt}
\appendixthree
\section*{Appendix 3}

Often the appendices contain technical details of the main results:
\begin{equation}
\label{C1}
a + b = c.
\end{equation}

\begin{remark}
This is a remark concerning equations~\eqref{A1} and \eqref{C1}.
\end{remark}

\begin{lemma}\label{facstrong}
The conditional density model $\mathcal{M}$ of $\S$\,3 is sequentially strongly convex with $H_{k}(p)\left( z\right)
\equiv p\left( a_{k}\mid \overline{l}_{k},\overline{a}_{k-1}\right) $.\vspace*{-10pt}
\end{lemma}

\bibliographystyle{biometrika}
\bibliography{paper-ref}

 




\end{document}\vspace{0.4cm} 
